\subsection{A Right to Consume and a Right to Produce}
\label{sec:right-consume}
The competing design is a basic income, and the choice between them is usually argued on forecasts about how much work machines will take. It does not have to be, and the forecast is not the reason to prefer one.

A jobs guarantee is a public option: a thing the state provides directly, to anyone, at a standard it is answerable for, sitting in the same market as the private version and setting the standard that version has to beat. A public school is one; a public health insurer is another. A basic income is not that. It is a transfer, and what it buys is whatever is for sale. Consuming public resources is not the same as being public — a charter school runs on public money and is not a public school — and a transfer spent at a private employer's store does not make the store public.

What employment carries besides money is attachment. A person in a guaranteed job is tied to other people by the work: colleagues, a schedule, a place, something produced that somebody needed — the last of which holds where the placement is real work and not make-work, a condition on the program rather than a property of the form. A transfer supplies the purchasing power and none of the rest.

A transfer also leaves a person outside the one body of law that reaches the terms on which they are dealt with. Anti-discrimination duties attach to employment. They do not attach to income. A guaranteed job puts somebody inside that architecture; a payment leaves them in a market, dealing with whoever will sell to them, on whatever terms. The architecture is contested and full of holes — it reaches employers above a size, and what it demands of a seller has always been narrower than what it demands of an employer — but a hole in a protection is something to litigate. An income entitlement supplies a claim to eligibility and payment, and a misapplied rule can be litigated on that claim. What it does not supply is the employment relationship, which is where the anti-discrimination duties attach, and that is the difference between the two arrangements this section rests on.

Both are conditional, at two points each: who can end the funding, and who can deny one person access. At the first point the two fail alike. A budget line can be withdrawn from everybody at once by an ordinary majority, a constitutional income would be harder to withdraw than a statutory guarantee, and a guarantee's funding fails the same way, which is why the previous section required it to be unconditional and centrally guaranteed. At the second point they differ in kind. Nobody denies one person a basic income; eligibility is a rule, and where it is misapplied the remedy is a payment claim. A placement is refused to a person, by somebody, on a date, which exposes an official and leaves a record for an appeal to attach to. That is a lever, and a dangerous one under a hostile administration, and it is also the only one of the two designs that leaves a decision-maker behind.

A decision someone can contest gives dependence a point of resistance. On that comparison, the advantage belongs to the arrangement that leaves a usable record and someone who must answer for it. It is also, for the same reason, the more demanding proposal to build: everything in the comparison rests on the appeal being real, and the previous section says what makes it real.

That standard — a record, a decision-maker, an appeal — has a court behind it. When the Supreme Court kept a Federal Reserve governor in her post in 2026 against a removal for cause (§\ref{sec:build-costs}), what it held she was owed was not the post but the artifact: “some explanation of the evidence at issue, some avenue for a response, and a deadline by which a response would be due” \autocite{cook2026}. Half of that case does not travel: the Fed's protection rests on a pedigree no body chartered later can acquire. The other half does. What an adverse decision about one person's tenure has to produce before it counts is not tied to that pedigree, and a placement refused is an adverse decision about one person's access to a public post.

The argument concerns a right to produce, which a transfer does not supply and a guarantee does. A basic income is the right to consume, and a right to consume is exercised in whatever market the concentration described at §\ref{sec:concentration-power} leaves standing.

I put a jobs guarantee ahead of a basic income on those grounds. One remaining measure does work neither design can: broadened social insurance — unemployment benefits, healthcare, housing — reaches the people who cannot work or choose not to, whom a guarantee by construction does not reach at all. That is the standard case for an income, it is not answered here, and a guarantee needs the insurance beside it.
